\scanpage{28}\mspage{27}
\uline{Cont} 3 Si $F$ est une \deleted{figille} figure, et $F'$ une partie fermée de l'ens ordonné $F$ (ie. $X,Y\in F$, $Y\in F'$, $\deleted{X}\leq Y\implies X\in F$) alors $F'$ est une figure.

On dit que $F'$ est une sous-figure -- ainsi, les sous figures de $F$ sont en corr. biunivoque avec les parties fermées de l'ens ordonné $F$.

Attention, si $F',F$ sont deux figures, on peut avoir $F'\subset F$ sans que $F'$ soit une sous figure. (Dans le cas de l'exemple-type, prendre $F$ formé de $X_0\subset X_1$, $X_1$ une variété connexe, $F'=\{X_1\}$.)

% slanted
\margin{ici, $F'\subset F$ mais $F'\nleq F$ et $F',F$ ne sont $\ne$ compatibles, au sens plus bas. --}

On écrira \deleted{$\unsure{\leq|\subset}$} $F'\leq F$ pour $F'$ sous-figure de $F$. C'est une relation d'ordre entre figures.

\deleted{Soient $F,\unsure{C|G}$ deux figures. \deleted{\unsure{(plus gén. la défin.|(plus gén. la définition} d'équivalence} \inserted{On dit que} $F$ et $G$ sont \deleted{\unsure{équivalentes}} \inserted{ens\textsuperscript{t}.} \inserted{compatibles}, si $F\cup G$ est}
